\section{Literature Review}\label{sec:literature}

\subsection{Catastrophic Health Expenditure in Developing Countries}

The conceptual framework for catastrophic health expenditure originates with \citet{Wagstaff2003} and \citet{Xu2003}, who define CHE as OOP spending exceeding a fixed share of household resources---total consumption or capacity-to-pay (consumption net of subsistence food expenditure). This binary threshold approach has been applied extensively across low- and middle-income countries, generating a large body of evidence on CHE prevalence and its correlates \citep{Wagstaff2008, WHO2019}. Common findings include that CHE disproportionately affects the uninsured, the elderly, households with chronic illness, and those in lower income quintiles.

The binary threshold approach has well-known limitations. \citet{Flores2008} demonstrate that CHE prevalence is sensitive to the choice of threshold and denominator, while \citet{Moreno2011} argue that mean-based regression on the binary indicator discards information about spending intensity. Households spending 11\% and 90\% of consumption on health are treated identically if the threshold is 10\%, yet their financial distress differs qualitatively.

\subsection{Quantile Regression in Health Expenditure Analysis}

Quantile regression, introduced by \citet{Koenker1978} and formalized in \citet{Koenker2005}, estimates the conditional quantile function, allowing covariate effects to vary across the distribution. In health expenditure analysis, this is particularly valuable because the OOP distribution is typically zero-inflated, right-skewed, and heteroskedastic.

\citet{Jones2013} survey quantile regression methods for health economics applications, emphasizing the distinction between conditional and unconditional quantile effects. \citet{Machado2019} develop a quantile regression framework for count data with applications to health care utilization. In the developing-country context, \citet{Bredenkamp2011} use quantile regression to analyze OOP spending in the Western Balkans, documenting that insurance effects are concentrated in the upper tail.

A critical methodological challenge arises when the outcome has a mass point at zero. Standard Koenker--Bassett quantile regression is degenerate at quantiles below the zero-mass fraction, producing uninformative estimates \citep{MachadoSilva2005}. Two approaches address this. First, censored quantile regression \citep{Powell1986, Chernozhukov2002} treats zero as a censoring point, producing consistent estimates of the latent quantile function. Second, the two-part model \citep{Manning1987, Mullahy1998} separates the extensive margin (participation) from the intensive margin (spending conditional on participation), allowing different covariates to drive each decision. We adopt the two-part framework as a complement to our conditional quantile regression, restricting interpretation of CQR results to quantiles above the zero-mass threshold.

The recentered influence function (RIF) regression of \citet{Firpo2009} estimates unconditional quantile effects---how a covariate shifts the population-level quantile. While conditional quantile regression answers ``at what point in the conditional distribution does a covariate matter?'', RIF regression answers the population-level question implied by many policy discussions. We note this distinction but focus on conditional quantile regression, as the insurance--quintile interaction effects we estimate are inherently conditional objects. A RIF extension would provide the population-level complement and is a natural direction for future work.

\subsection{Health Financing in Peru}

Peru's health system is characterized by segmentation. SIS, created in 2002, provides subsidized coverage to households in poverty and extreme poverty, covering primary care and selected hospital services at public facilities \citep{Petrera2013}. The legal framework for universal coverage was established by Law 29344 (Ley Marco de Aseguramiento Universal en Salud, 2009), which created the Plan Esencial de Aseguramiento en Salud (PEAS) as the minimum benefit package. EsSalud, the social security scheme, covers formal-sector workers and their dependents, providing more comprehensive benefits. Private insurance (EPS) covers a small fraction of the population, primarily in Lima.

\citet{EspinozaPajuelo2024} analyze informal payments in Peruvian health facilities using the 2018 ENAHO, finding that SIS beneficiaries incur undisclosed costs for nominally free services. \citet{Petrera2013} document persistent OOP spending among SIS enrollees, particularly for medications and diagnostic services not covered by PEAS. In the Latin American context, \citet{Knaul2006} study Mexico's Seguro Popular---a subsidized insurance scheme targeting the poor, analogous to SIS---finding that coverage expansion reduced CHE prevalence but with heterogeneous effects across the income distribution.

A critical institutional feature is SIS targeting imprecision. The SISFOH system uses proxy means testing to determine eligibility, but substantial leakage---enrollment of non-poor households---has been documented \citep{Cortez2008, Francke2013}. This leakage complicates the interpretation of SIS effects at upper consumption quintiles: positive SIS enrollment at Q4--Q5 may reflect targeting errors rather than legitimate enrollment, and any measured financial protection at these quintiles conflates the insurance effect with the composition of the (mis)targeted population.

\subsection{Identification Challenges with Observational Insurance Data}

Estimating the effect of health insurance on OOP spending faces fundamental identification challenges. Insurance enrollment is non-random: means-tested programs like SIS select on poverty status, while contributory schemes like EsSalud select on formal employment. \citet{Wagstaff2010} formalizes these challenges for observational studies and surveys quasi-experimental approaches (RDD at eligibility thresholds, enrollment mandates) that can identify causal effects under specific assumptions. \citet{Acharya2012} provide a systematic review of health insurance impact evaluations in low- and middle-income countries, finding that observational studies consistently overestimate protective effects due to uncontrolled selection.

Our paper is explicitly descriptive. We do not claim to identify causal effects of insurance on OOP spending. Instead, we document conditional distributional patterns that characterize where financial protection gaps are concentrated---information that binary CHE indicators and mean-based regressions cannot deliver, but that stops short of causal attribution.
